\section{Cross-Prime Logarithmic Coherence}

\subsection{The signed input exponent signal}
Let
\[
 S_{ab}=\{q:q\mid ab\}.
\]
Define
\[
 \nu_q=v_q(b)-v_q(a).
\]
Then
\[
 -\frac ba=-\prod_{q\mid ab}q^{\nu_q}.
\]
For every output prime $p\mid c$, the same input exponent packet produces the deep principal unit
\[
 1-\frac ca.
\]
The output persistence is therefore a simultaneous cross-prime cancellation phenomenon.

\subsection{Odd logarithmic rows}
For odd $p^e\parallel c$, define the finite-depth row coefficient
\[
 \ell_{p,e}(q)
 =
 \frac1p\log_p\langle q\rangle_p
 \pmod{p^{e-1}}.
\]
The balanced relation gives
\[
 \boxed{
 \sum_{q\mid ab}
 \nu_q\ell_{p,e}(q)
 \equiv0
 \pmod{p^{e-1}}.
 }
\]
Define the truncated valuation on $\Z/p^{e-1}\Z$ by
\[
 \bar v_p(0)=e-1,
 \qquad
 \bar v_p(x)=\min\{e-1,v_p(\widetilde x)\}
\]
for any integer lift $\widetilde x$ of a nonzero residue. Then
\[
 \boxed{
 \bar v_p(\ell_{p,e}(q))
 =
 \min\{e-1,v_p(q^{p-1}-1)-1\}.
 }
\]
Exceptional entrywise depth is a generalized Wieferich event.

\subsection{Weighted dyadic row}
For $2^e\parallel c$ with $e\ge2$ and an odd input prime $q$, define
\[
 \ell_{2,e}(q)
 =
 \frac14
 \log_2\left((-1)^{(q-1)/2}q\right)
 \pmod{2^{e-2}}.
\]
With the analogous truncated valuation on $\Z/2^{e-2}\Z$,
\[
 \bar v_2(\ell_{2,e}(q))
 =
 \min\{e-2,v_2(q^2-1)-3\}.
\]
At $e=1$, the output contact depth is zero and this row is omitted. A separate sign packet supplies the missing first dyadic layer. This is the weighted dyadic branch in its final logarithmic form.

\subsection{Sequential factorization of the powerful part}
The factorization is defined prime by prime. For an odd repeated output prime $p^e\parallel c$, set
\[
 d_p=e-1.
\]
For a repeated dyadic prime $2^e\parallel c$, remove one sign layer and set
\[
 s_2=1,
 \qquad
 d_2=e-2.
\]
If $2\nmid M$, take $s_2=0$. Put $\epsilon_2=2^{s_2}$. Let $K_{\mathrm{Kum}}$ be the kernel size of the balanced Kummer class and put
\[
 r_p=\min\{d_p,v_p(K_{\mathrm{Kum}})\}.
\]
The balanced order formula implies
\[
 K_{\mathrm{Kum}}\mid M/\epsilon_2,
\]
so $K_{\mathrm{Kum}}=\prod_{p\mid c}p^{r_p}$. Write $K=K_{\mathrm{Kum}}$. Divide the input exponent vector by $K$. Define $a_{p,q}$ as the original odd or weighted dyadic coefficient, now reduced modulo $p^{d_p-r_p}$. Dividing the relation vector reduces the needed relation modulus. It does not divide the coefficient row by $K$. Define
\[
 w_p
 =
 \min\left
 \{
 d_p-r_p,
 \min_{q\mid ab}\bar v_p(a_{p,q})
 \right\},
 \qquad
 \chi_p=d_p-r_p-w_p.
\]
The normalized global factors are
\[
 K=\prod_{p\mid c}p^{r_p},
 \qquad
 W=\prod_{p\mid c}p^{w_p},
 \qquad
 C=\prod_{p\mid c}p^{\chi_p}.
\]
The balanced Kummer order formula guarantees that the dyadic sign layer is removed before the principal $2$-adic Kummer depth, so these exponents partition the contact depth. Hence
\[
 \boxed{
 M=\epsilon_2KWC.
 }
\]
The order of removal matters. The local logarithmic defect is measured only after the Kummer preload has been divided out. The resulting row is then divided by its common logarithmic depth, leaving a primitive row modulo $p^{\chi_p}$.

\subsection{The coherence character}
Set $\chi_p=0$ for absent or squarefree output channels. When the residual modulus is one, use the trivial row and character. After normalization, every output prime with nontrivial residual depth produces a primitive row
\[
 g_p
 \in
 (\Z/p^{\chi_p}\Z)^{S_{ab}}
\]
satisfying
\[
 g_p\cdot\nu_0
 \equiv0
 \pmod{p^{\chi_p}},
\]
where $\nu_0$ is the Kummer-normalized input exponent vector. The Chinese remainder theorem combines the local rows into one primitive row
\[
 g\in(\Z/C\Z)^{S_{ab}}.
\]
Define the coherence character
\[
 \boxed{
 \chi_{\mathrm{coh}}(x)
 =
 \exp\left(\frac{2\pi i}{C}g\cdot x\right).
 }
\]
Its kernel is the adelic coherence lattice
\[
 \boxed{
 \Lambda_{\mathrm{coh}}
 =
 \{x\in\Z^{S_{ab}}:g\cdot x\equiv0\pmod C\}.
 }
\]
Since $g$ is primitive modulo $C$,
\[
 [\Z^{S_{ab}}:\Lambda_{\mathrm{coh}}]=C
\]
and
\[
 \Z^{S_{ab}}/\Lambda_{\mathrm{coh}}\cong\Z/C\Z.
\]
The normalized exponent vector lies in the lattice.

\subsection{Worked lattice for \texorpdfstring{$3+125=128$}{3+125=128}}
Here
\[
 M=64
\]
and
\[
 (\epsilon_2,K,W,C)=(2,1,1,32).
\]
The input vector is
\[
 \nu_0=(-1,3)
\]
on the prime coordinates $(3,5)$. The residual row is
\[
 (29,31)\pmod{32}.
\]
Thus
\[
 \boxed{
 \Lambda_{\mathrm{coh}}
 =
 \{(x_3,x_5)\in\Z^2:29x_3+31x_5\equiv0\pmod{32}\}.
 }
\]
Indeed,
\[
 29(-1)+31(3)=64\equiv0\pmod{32}.
\]
The five residual dyadic contact layers are exactly the index of this lattice.

\subsection{Generalized Wieferich versus collective coherence}
There are two qualitatively different sources of large output depth.

\begin{enumerate}[leftmargin=1.5em]
\item An entrywise defect can make every coefficient in one row divisible by a large power of $p$. This is the factor $W$.
\item Ordinary coefficients can cancel coherently against the specific exponent vector. This is the factor $C$.
\end{enumerate}

In the recovered 361-row coherence corpus, the recorded $W$ factor is usually one while $C$ can be large. The corpus is selected, rather than a statistical sample of all primitive triples. Its frequencies do not establish a generic population law. For its $W=1$ rows, the residual depth is represented by collective cancellation rather than an entrywise Wieferich defect.

\subsection{Why abstract lattice geometry is insufficient}
Let $v\in\Z^s$ be primitive, with $s\ge2$. For every $C\ge2$, there exists a primitive row $g\pmod C$ satisfying
\[
 g\cdot v\equiv0\pmod C.
\]
Thus a fixed short vector can lie in primitive congruence lattices of arbitrarily large index.

\begin{theorem}[Arbitrary-index annihilator]
Abstract lattice geometry cannot bound the coherence index from the length of the known relation vector alone.
\end{theorem}

The proof completes $v$ to a unimodular basis and chooses a primitive row annihilating its first basis coordinate modulo $C$.

The missing theorem must use the arithmetic origin of $g$: its entries are simultaneous $p$-adic logarithms of the actual input primes, and the output primes are exactly the divisors of $a+b$.

\subsection{Why positive divisibility is circular}
If positive integers $A$ and $B$ satisfy $c\mid A+B$, then
\[
 A+B\ge c,
\]
so
\[
 \max(A,B)\ge c/2.
\]
A direct search for a small positive $S$-unit relation divisible by $c$ cannot produce an $abc$ saving. Signed cancellation, a determinant, a regulator, or a functional equation is necessary.

\subsection{Exact reformulation of strong \texorpdfstring{$abc$}{abc}}
Since
\[
 c=R_c\epsilon_2KWC
\]
and $\epsilon_2\le2$ while $K\le\log_2c$, the all-$\eps$ strong $abc$ statement is equivalent, after the usual absorption of bounded and logarithmic factors into an arbitrarily small power, to the target
\[
 \boxed{
 WC
 \le
 C_\eps R_aR_bR^\eps.
 }
\]
The logarithmic absorption is proved explicitly in Section~\ref{sec:bridge}, without assuming an initial relation $c\ll R$. Conversely, $WC\le M_c$ makes the reverse implication immediate. In the recorded examples with $W=1$, this leaves the coherence conductor bound
\[
 C\le C_\eps R_aR_bR^\eps.
\]

\subsection{Adelic coherence regulator object}
Package the surviving data as
\[
 \boxed{
 \mathfrak R_{a,b,c}
 =
 \left(
 U_{ab},
 \beta^{1/K},
 \{\chi_p\}_{p\mid c},
 \chi_{\mathrm{coh}},
 \Lambda_{\mathrm{coh}},
 \norm{\cdot}_\infty
 \right).
 }
\]
Here $U_{ab}$ is the input $S$-unit group, the local characters are the normalized logarithmic rows, and the archimedean norm records input height.

The open theorem is an adelic slope or regulator inequality bounding the finite conductor $\log C$ by the input radical and an arbitrarily small global leakage term. General simultaneous linear-form theory contains the right kinds of local and archimedean quantities, but known general estimates are much too coarse for the $R^\eps$ scale.

